% Módulo científico sin preámbulo ni portada autónoma.
% Dependencia causal: atlas, operadores K_D/K_\Omega y contraste metrológico.

\chapter{Equivariant realization of species and mass operators: genealogical fibres, spectral measures, complex poles, and material compositions}
\label{hmtmasscl:chap-realizacion-equiv}

\section{Internal domain of the material realization and classification of its codomains}

\subsection{Separation among spectral line, multiplet, distribution, complex pole, categorical object, and criterion of impossibility}

Let \(\mathcal P\) be the set of individualized physical states. An elementary
state is represented by the descriptor
\[
 X=(\mathfrak c,\rho,g,f,q,\epsilon,\varsigma,\mathfrak a),
\]
which preserves representation class \(\mathfrak c\), internal representation
\(\rho\), generation \(g\), flavor \(f\), charge \(q\), sheet orientation
\(\epsilon\), statistical sector \(\varsigma\), and coupling data
\(\mathfrak a\). A composite state adds ordered constituents, composition
tree, common boundary, and collective quantum numbers. A resonance
adds the open channels and the causal prescription for the resolvent.

The physical realization must address three distinct questions:

\begin{enumerate}
\item which fibre, orbit, or multisection of the atlas corresponds to the descriptor;
\item which positive operator publishes the mass on that fiber;
\item which coupling produces a scalar, a multiplet, a distribution, or
      a complex pole.
\end{enumerate}

Confusing these questions introduces a retrospective selector. The correct
answer is obtained through a natural relation, a fiber operator and a
spectral compression, in that order.

\subsection{Precedence of the genealogical atlas with respect to any metrological identification}

The General Mass Law does not start from a table of names or a list of
experimental values. Its internal domain is the finite chain
\[
 \mathcal R_{324}\longrightarrow \mathcal C_{81}
 \longrightarrow \mathcal F_{56}\longrightarrow \mathcal M_{13},
\]
where \(\mathcal R_{324}\) contains the \(4\) orientations of each of the
\(81\) cells, \(\mathcal F_{56}\) groups equivalent route histories, and
\(\mathcal M_{13}\) preserves the operator multisections. The null sector
bifurcates before the positive character is applied. Therefore, neither an
abbreviated historical table nor a taxonomy of physical names can replace this
domain.

\section{Fiber product between species descriptors and routes of the 81/56/13/324 atlas}

\subsection{Exact non-voidness condition and prospective selection without illicit point selector}

Let \(\mathcal D\) be the finite space of internal data already constructed before the
contrast: representation, orientation, memory, boundary, charge, flavor,
color, generation, and statistical sector. Denote by
\[
 d_{\mathcal P}:\mathcal P\longrightarrow\mathcal D,
 \qquad
 d_{\mathcal R}:\mathcal R_{324}\longrightarrow\mathcal D
\]
the physical descriptor maps and the genealogical descriptor. The object
of realization is defined by the fiber product
\[
 \mathcal J
 =\mathcal P\times_{\mathcal D}\mathcal R_{324}
 =\{(X,\gamma):d_{\mathcal P}(X)=d_{\mathcal R}(\gamma)\}.
\]
The fiber of a species is
\[
\mathscr S(X)=\{\gamma\in\mathcal R_{324}:(X,\gamma)\in\mathcal J\}.
\]
A descriptor (X) is called realizable by the atlas when
\[
 d_{\mathcal P}(X)\in\operatorname{im}d_{\mathcal R}.
\]
This condition is equivalent to \(\mathscr S(X)\ne\varnothing\). It does not follow from merely writing the fiber product: it constitutes the exact existence condition that must be verified for each physical descriptor. The value of \(\mathscr S\) may be a route, an orbit, or a multisection. A point section is not forced when the internal symmetry requires preservation of a multiplicity.

\subsection{Equivariance of fibers, orbits, and multisections under refinement}

\begin{definition}[Admissible realization relation]
The relation \(\mathscr S\) is admissible when it simultaneously satisfies:
\begin{enumerate}
\item depends solely on data constructed before the external table is opened;
\item commutes with particle--antiparticle conjugation;
\item is equivariant under cellular inversion and the internal
      gauge actions;
\item it is compatible with refinement, nonadic return and dodecaphasic
      reading;
\item preserves every multiplicity not eliminated by a physical coupling;
\item remains invariant under permuting or replacing the external values of
      mass and width.
\end{enumerate}
\end{definition}

\begin{theorem}[Prospective realization by fibered product]
If \(d_{\mathcal P}\) and \(d_{\mathcal R}\) are equivariant with respect to a symmetry \(G\) and do not depend on external magnitudes, then \(\mathcal J\) is a \(G\)-subobject of \(\mathcal P\times\mathcal R_{324}\). For every realizable descriptor \(X\), the nonempty fiber \(\mathscr S(X)\) is \(G_X\)-stable, and the assignment \(X\mapsto\mathscr S(X)\) is independent of the result subsequently compared.
\end{theorem}

\begin{proof}
For \((X,\gamma)\in\mathcal J\) one has
\(d_{\mathcal P}(X)=d_{\mathcal R}(\gamma)\). Equivariance gives
\[
 d_{\mathcal P}(gX)=g\,d_{\mathcal P}(X)
 =g\,d_{\mathcal R}(\gamma)=d_{\mathcal R}(g\gamma),
\]
and, consequently, \((gX,g\gamma)\in\mathcal J\). If \(g\in G_X\), the
second component remains in \(\mathscr S(X)\). Independence of the
external values is immediate because none of them belongs to the domain
of the two maps of the fiber product.
\end{proof}

\section{Intrinsic spectral classification of material species}

\subsection{Pointwise support decomposition, irreducible block and fiber image measure}

Let \(\mathcal A_{\rm rut}\) be the finite commutative algebra generated by the
internal route observables: signature \(\mathbb Z^5\), character, chronology,
\(K_D\), \(K_\Omega\), orientation, memory, and compatible readers. In a
fiber \(F\), one defines
\[
 \gamma\sim_{\rm sp}\eta
 \quad\Longleftrightarrow\quad
 a(\gamma)=a(\eta)\quad\text{for every }a\in\mathcal A_{\rm rut}.
\]
For each class \(\lambda\in F/\!\sim_{\rm sp}\),
\[
 P_\lambda=\sum_{\gamma\in\lambda}
 |\gamma\rangle\langle\gamma|
\]
is a spectral projector; the \(P_\lambda\) are orthogonal and
\(\sum_\lambda P_\lambda=I_F\).

\begin{theorem}[Maximum internal individualization]
Every classifier constructed exclusively from observables of \(\mathcal A_{\rm rut}\) factors uniquely through the quotient \(F/\!\sim_{\rm sp}\). Consequently, two descriptors with distinct complete signatures belong to distinguishable sectors; two descriptors that agree for all generators of the algebra cannot be separated by a rule constructed from those same observables.
\end{theorem}

\begin{proof}
The characters of the finite commutative algebra separate exactly its
joint spectrum. A functional classifier of its generators is constant
on each fiber of the joint spectral map and therefore factors through the
quotient. If two routes have the same character on the entire algebra, no
element of it can separate them.
\end{proof}

\subsection{Distinction between scalar mass, multiplet, and spectral distribution}

This theorem closes the logical limit of the classification. When a
physical representation requires two states belonging to the same
spectral class to be distinguished, it must exhibit an additional observable or a coupling that
enlarges \(\mathcal A_{\rm rut}\). The absence of such an observable does not authorize
the external mass to be used as a selecting label.

\section{Natural refinement of the enriched record in the positive towers \texorpdfstring{\(90/120\)}{90/120}}

\subsection{Positive enumeration of the semigroup \texorpdfstring{\(\langle90,120\rangle\)}{<90,120>} and exclusion of mode 0}

The existence of the semigroup \(\langle90,120\rangle\) is not sufficient by itself:
the map assigning coefficients to an enriched route must be constructed. Let
\[
 \mathcal H_{90,120}
 =\{90a+120b:a,b\in\mathbb N_0,\ (a,b)\ne(0,0)\}
 =\{90,120\}\cup\{30r:r\geq6\}
 =\{h_0<h_1<h_2<\cdots\}
\]
its strictly increasing enumeration, where \(\mathbb N_0=\{0,1,2,\ldots\}\),
and let \(s_{h_n}>0\) be the declared normalization of the corresponding mode. The
compatible record of a route
\(\Gamma\) provides a sequence of finite integer data
\[
 \mathcal L(\Gamma)=(\ell_0(\Gamma),\ell_1(\Gamma),\ldots),
\]
in which each \(\ell_n\) brings together return, carry, winding, memory and boundary
at refinement \(n\). To fix the encoding unambiguously, let
\[
 \nu(z)=
 \begin{cases}2z,&z\geq0,\\-2z-1,&z<0,\end{cases}
 \qquad
 \pi_{\rm C}(u,v)=\frac{(u+v)(u+v+1)}2+v,
\]
and define \(\iota:\mathbb Z^d\to\mathbb N_0\) by recursive application of
\(\pi_{\rm C}\) to \(\nu(z_1),\ldots,\nu(z_d)\). Then
\[
 q_n(\Gamma)
 =\frac{\iota(\ell_n(\Gamma))}
 {1+|\iota(\ell_n(\Gamma))|}\in(-1,1).
\]
The minimal realization of tower is
\[
 \eta_{h_n}(\Gamma)
 =\frac{3^{-(n+1)}q_n(\Gamma)}{s_{h_n}},
 \qquad
 \mathfrak T(\Gamma)
 =(\eta_h(\Gamma))_{h\in\mathcal H_{90,120}}.
\]

\subsection{Absolute convergence, prefix naturality, and effective tail bound}

\begin{theorem}[Convergence, naturality and metrological blindness]
For every enriched route,
\[
 \sum_{h\in\mathcal H_{90,120}}|\eta_h(\Gamma)s_h|
 \leq\sum_{n\ge0}3^{-(n+1)}=\frac12.
\]
If the refinement prolongs the record by prefixes, \(\mathfrak T\) is natural
with respect to the truncations. Moreover, no external mass value intervenes
in \(\mathfrak T\), and the tail after \(N\) modes satisfies
\[
 \sum_{n>N}|\eta_{h_n}s_{h_n}|
 \leq\frac{3^{-(N+1)}}2.
\]
\end{theorem}

\begin{proof}
We have \(|q_n|<1\); the two bounds are geometric sums. Prefix
compatibility preserves all coefficients already constructed and adds those of the
new level, which yields the naturality square with the truncation of
\(\mathcal E_{90,120}\). The formula contains only the record, the explicit
coding, and the declared modal normalization.
\end{proof}

The preceding construction fixes an explicit, convergent minimal realization relative to the encoding \(\iota\) and the declared normalizations \(s_h\). A dynamics that modifies its weights must declare the operator effecting that modification and prove naturality and convergence anew; it cannot reconstruct the weights from an experimental residual.

\section{Bidirectional operators \texorpdfstring{\(K_D\) and \(K_\Omega\)}{K_D and K_Omega} and joint spectral measure}

\subsection{PVM joint, compression to POVM and scalar distribution induced by a state}

For \(F_X=\mathscr S(X)\), let
\[
 \mathcal H_X=\ell^2(F_X)\otimes V_{\rho_X}
\]
the realization space, where \(V_{\rho_X}\) carries the corresponding internal
representation. Each route \(\gamma\in F_X\) has signature
\[
 \sigma(\gamma)
 =(n_A,s_C,k_{\Delta_4},b_{120},b_\zeta)\in\mathbb Z^5
\]
and refinement coefficients
\(\eta(\gamma)\in\mathcal E_{90,120}\). The basal operator is
\[
 \widehat M_X^{(0)}
 =\sum_{\gamma\in F_X}
 m_{\rm clk}\,
 \chi_{\rm full}(\sigma(\gamma),\eta(\gamma))
 |\gamma\rangle\langle\gamma|\otimes I_{V_{\rho_X}}.
\]
The scale \(m_{\rm clk}\) descends from the absolute branch
action--clock--mass, including the decade \(10^{-34}\); it is not extracted from an
experimental mass.

The readers \(K_D\) and \(K_\Omega\) induce exact diagonal operators
\(B_X^D\) and \(B_X^\Omega\). The two orientations publish
\[
 \widehat M_X^{\beta,r}
 =R_{\rm act}^{B_X^r/2}\widehat M_X^{(0)}
  R_{\rm act}^{B_X^r/2},
 \qquad r\in\{D,\Omega\}.
\]
Both are positive and commute in the route basis. By the joint spectral theorem
there exists a unique projective spectral measure
\(E_X^{D,\Omega}\) over \(\mathbb R_+^2\) such that
\[
 \widehat M_X^{\beta,D}
 =\int x\,dE_X^{D,\Omega}(x,y),
 \qquad
 \widehat M_X^{\beta,\Omega}
 =\int y\,dE_X^{D,\Omega}(x,y).
\]
At the base of routes,
\[
 E_X^{D,\Omega}(A)|\gamma\rangle
 =\mathbf 1_A\!\left(m_D(\gamma),m_\Omega(\gamma)\right)
  |\gamma\rangle.
\]
This is the canonical bidirectional output: it preserves the complete pair and does not
presuppose a scalar functional in advance.

For a channel \(P_{X,a}\), the compressed spectral measure
\[
 F_{X,a}^{D,\Omega}(A)
 =\left.P_{X,a}E_X^{D,\Omega}(A)P_{X,a}\right|_{P_{X,a}\mathcal H_X}
\]
is a POVM on the prepared subspace. For a density \(\rho_X\)
supported on that subspace, the joint scalar distribution is
\[
 \mu_{X,a}^{D,\Omega}(A)
 =\operatorname{Tr}\!\left(
 \rho_X F_{X,a}^{D,\Omega}(A)
 \right).
\]
Every declared functional \(f:\mathbb R_+^2\to\mathbb R_+\) induces
\(\mu_{X,a}^f=f_\#\mu_{X,a}^{D,\Omega}\) without modifying the source PVM.

\subsection{Mixed moments, point spectral support and conservation of multiplicities}

Let
\[
 \boldsymbol\beta_X(\gamma)
 =(m_D(\gamma),m_\Omega(\gamma)).
\]
The uniform image measure of the fiber is
\[
 \nu_X=(\boldsymbol\beta_X)_\#
 \left(\frac1{|F_X|}\sum_{\gamma\in F_X}\delta_\gamma\right)
 =\frac1{|F_X|}\sum_{\gamma\in F_X}
 \delta_{(m_D(\gamma),m_\Omega(\gamma))}.
\]
Its transform and its mixed moments are
\[
 Z_X(s,t)=\int e^{sx+ty}\,d\nu_X(x,y),
 \qquad
 \mathfrak m_{p,q}(X)=\int x^py^q\,d\nu_X(x,y).
\]
Since \(\nu_X\) has finite support, \(Z_X\), or equivalently the complete
family of mixed moments, determines the joint multiset of reader pairs
up to permutation of routes. The pushforward measure is therefore a
complete invariant of the operator pair; it is not an average of masses.

If \(u:F_X\to F_Y\) is a bijection preserving record, orientation, and
both readers, the unitary \(U_u|\gamma\rangle=|u\gamma\rangle\) satisfies
\[
 U_uE_X^{D,\Omega}(A)=E_Y^{D,\Omega}(A)U_u,
 \qquad
 \nu_X=\nu_Y.
\]
The same equality holds for the physical distributions when density and
projector are transported by \(U_u\). This naturality is asserted only for
the intertwiners actually constructed; a coincidence of
cardinalities does not by itself create an equivalence of fibers.

\section{Symmetry Scalarization and Isotypic Decomposition of the Fibers}

\subsection{Conditional symmetry expectation and criterion for reduction to a spectral line}

Let \(G_X\) be the compact group of unbroken symmetry of the fiber and
\(U_X:G_X\to\mathcal U(\mathcal H_X)\) its representation. The conditional
symmetry expectation is defined by
\[
 \mathbb E_X(T)
 =\int_{G_X}U_X(g)TU_X(g)^*\,d\mu_X(g),
\]
where \(\mu_X\) is the normalized Haar measure. The physical mass operator
before opening channels is defined, when the state type declares a
functional \(\Phi_X:\mathbb R_+^2\to\mathbb R_+\), by
\[
 \mathcal M_X
 =\mathbb E_X\!\left(
   \Phi_X(\widehat M_X^{\beta,D},\widehat M_X^{\beta,\Omega})
  \right).
\]
The functional calculus is carried out with the joint measure
\(E_X^{D,\Omega}\). In a representation containing an intertwiner that
exchanges readers, one may take \(\Phi_X(x,y)=\sqrt{xy}\). If the physical
type does not declare \(\Phi_X\), the complete output remains the joint PVM;
no scalar is fabricated.

\begin{theorem}[Scalar closure, multiplet, and distribution]
The operator \(\mathcal M_X\) belongs to the commutant
\(U_X(G_X)'\). If
\[
 \mathcal H_X\simeq
 \bigoplus_\lambda V_\lambda\otimes\mathbb C^{m_\lambda},
\]
then
\[
 \mathcal M_X=
 \bigoplus_\lambda I_{V_\lambda}\otimes M_{X,\lambda}.
\]
In an irreducible sector of multiplicity \(1\), the restriction is a scalar.
In a reducible sector, the exact output is the spectrum of the blocks
\(M_{X,\lambda}\), with their multiplicities. There is no indeterminacy:
there is a distinct codomain.
\end{theorem}

\begin{proof}
Invariance of the Haar measure implies
\(U_X(h)\mathbb E_X(T)U_X(h)^*=\mathbb E_X(T)\) for every \(h\in G_X\).
The decomposition follows from the complete reducibility theorem and Schur's
lemma. The irreducible case produces a multiple of the identity; the
reducible case preserves exactly the multiplicity blocks.
\end{proof}

Let \(P_{X,a}\) be the projector determined by the prepared state and the physical channel or
apparatus \(a\). The observable spectral measure is
\[
 \mu_{X,a}(B)=
 \langle\Psi_X,
 P_{X,a}E_{\mathcal M_X}(B)P_{X,a}\Psi_X\rangle.
\]
There is a scalar line precisely when the compressed spectral support is
pointwise:
\[
 E_{\mathcal M_X}(\{m_X\})P_{X,a}=P_{X,a},
 \qquad\text{equivalently}\qquad
 (\mathcal M_X-m_X I)\,P_{X,a}=0.
\]
The isolated equality
\(P_{X,a}\mathcal M_X P_{X,a}=m_X P_{X,a}\) only fixes an expectation if the
subspace is not reducing and is not enough to prove a spectral line.
If the compression preserves several eigenvalues, the output is a multiplet or a
distribution. Replacing it with a numerical average would destroy physical
and genealogical information.

\subsection{Geometric mean of the readers only under exchange invariance demonstrated}

When the physical representation requires symmetry under the exchange
\(D\leftrightarrow\Omega\), the minimal operatorial aggregate is the geometric
mean
\[
 \widehat M_X^{\rm sym}
 =\widehat M_X^{\beta,D}\#\widehat M_X^{\beta,\Omega},
 \qquad
 A\#B
 =A^{1/2}\bigl(A^{-1/2}BA^{-1/2}\bigr)^{1/2}A^{1/2},
\]
on the common positive support. In the diagonal case,
\[
 \widehat M_X^{\rm sym}
 =R_{\rm act}^{(B_X^D+B_X^\Omega)/4}
  \widehat M_X^{(0)}
  R_{\rm act}^{(B_X^D+B_X^\Omega)/4}.
\]
This formula is a consequence of the explicit exchange hypothesis; it does not
replace the joint measure when that symmetry has not been established.

\begin{theorem}[Characterization of the symmetric aggregate]
The geometric mean is positive, symmetric, congruence-covariant for
every invertible \(S\), and self-dual:
\[
 A\#B=B\#A,
 \qquad
 S^*(A\#B)S=(S^*AS)\#(S^*BS),
 \qquad
 (A\#B)^{-1}=A^{-1}\#B^{-1}.
\]
Moreover, \(A\#A=A\). Consequently, in sectors where there is a
physical intertwiner exchanging the readers, the aggregate introduces no
parameter and favors neither orientation.
\end{theorem}

\begin{proof}
The congruence identity is understood for invertible \(S\). The identities
follow from the functional calculus of the positive root of
\(A^{-1/2}BA^{-1/2}\). In the commutative case, they reduce to
\(A\#B=(AB)^{1/2}\).
\end{proof}

There is also a canonical scalar aggregate under distinct axioms. If
\(s_n\) is continuous, permutation-invariant, coordinatewise
multiplicative, and normalized by \(s_n(cI)=c\), then
\[
 s_n(T)=\det(T)^{1/n}.
\]
Indeed, upon passing to logarithms, continuity and additivity make the
functional linear in the logarithms of the eigenvalues; symmetry makes all
coefficients equal, and normalization fixes their value at \(1/n\). For the pair of
readers, exchange symmetry gives the invariant
\[
 S_{D,\Omega}
 =\sqrt{
 \operatorname{ndet}(T_D)\operatorname{ndet}(T_\Omega)},
 \qquad
 \operatorname{ndet}(T)=\det(T)^{1/n}.
\]
This result classifies the complete fiber under the stated axioms. It is not
automatically identified with a mass line: that identification requires
the physical coupling to adopt precisely that functional.

\section{Null sector and precedence of the zero-mass chart with respect to the positive character}

\subsection{Separation of the photonic \texorpdfstring{\(U(1)\)}{U(1)} and adjoint gluonic \texorpdfstring{\(SU(3)\)}{SU(3)} representations}

The photon belongs to the electromagnetic null chart and the gluons to the adjoint
null chart of color. Their different representations, polarizations, and
gauge restrictions remain visible even though both masses are null.

\subsection{Unreachability of the null sector as a limit of the positive mass branch}

Nullity is not obtained by taking a limit of the positive character. The domain
is first separated into
\[
 \mathcal H_{\rm phys}=\mathcal H_0\oplus\mathcal H_+,
\]
with
\[
 \mathcal M|_{\mathcal H_0}=0,
 \qquad
 \mathcal M|_{\mathcal H_+}>0.
\]

\section{Beta closure of the electron and charged lepton sectors}

\subsection{Uniqueness of the electron in cell \texorpdfstring{\((5,5)\)}{(5,5)} and preservation of the governing beta mass}

The electron constitutes the central exception. Cell inversion has a unique fixed point, the cell \((5,5)\), and its physical fiber therefore has rank one. The remaining stable multisections do not, in general, admit an equivariant point selector. This difference is not a deficiency of the atlas: it is the mathematical distinction between a central section and a species realized as an orbit or multisection.

For the rank-one electronic fiber, the joint measure is punctual, the two
readers share the first component \(80\), the compression is scalar, and the beta closure is
recovered
\[
 m_e^\beta
 =0.5109989506900008086082571648\ldots\ {\rm MeV}.
\]
The base value remains as a construction stage; it does not replace the beta closure.

\subsection{Realizations of the muon and tau via charged sectors and spectral criterion of scalarity}

In the charged leptonic sector, the descriptors distinguish the central fiber,
the depth-\(G_2\) multisection, and the depth-
\(G_4\) multisection. The first produces the unique electronic line; the other two
produce the operators associated with the muonic and tauonic realizations. If
the coupling preserves more than one block, the correct prediction is its
spectrum; the physical label does not authorize choosing one of its eigenvalues by
numerical proximity.

\section{Chromatic algebra \texorpdfstring{\(SU(3)\)}{SU(3)}, quark flavours, and running-mass transport}

\subsection{Weyl qutrit algebra, Reynolds reduction, and singlet projectors}

The chromatic orbit is realized in the color qutrit. The Weyl operators
generate \(\mathfrak{sl}_3(\mathbb C)\); the traceless anti-Hermitian selection
yields \(\mathfrak{su}(3)\), and the finite connection transports the three
chromatic charts. The flavor mass must be invariant under this action.

Let \(U_{\rm col}:SU(3)\to\mathcal U(V_{\rm col})\). The chromatic reduction
is the Reynolds expectation
\[
 \mathbb E_{\rm col}(T)
 =\int_{SU(3)}U_{\rm col}(g)T U_{\rm col}(g)^*\,dg.
\]
In the fundamental irreducible representation,
\[
 \mathbb E_{\rm col}(\mathcal M_q)
 =m_q(\mu,\mathscr R)I_3.
\]
The published mass is therefore independent of red, green and blue. The
coordinates \(\mu\) and \(\mathscr R\) identify the scale and renormalization
scheme; they are not hidden ambiguities and must always accompany the
figure.

For mesons and baryons, the singlet projectors are
\[
 P_{q\bar q}^{(1)}
 =\frac1{3}|\Omega\rangle\langle\Omega|,
 \qquad
 |\Omega\rangle=\sum_{a=1}^{3}|a\bar a\rangle,
\]
and
\[
 P_{qqq}^{(1)}
 =|\varepsilon\rangle\langle\varepsilon|,
 \qquad
 |\varepsilon\rangle
 =\frac1{\sqrt6}\sum_{a,b,c}\varepsilon_{abc}|abc\rangle.
\]
Both eliminate the observable chromatic label without erasing the internal orbit
that intervened in the construction.

\subsection{Renormalization Connection and Ordered Transport Between Scale Fibers}

The dependence on scale is a section of a fibration over
\(t=\log\mu\). If \(\gamma_m(g(t))\) is the endomorphism induced by the
chromatic connection, transport is defined by
\[
 \nabla_{\rm RG}m_q
 =\left(\frac{d}{dt}+\gamma_m(g(t))\right)m_q=0,
\]
\[
 m_q(t_2)
 =\mathcal P\exp\!\left(-\int_{t_1}^{t_2}
 \gamma_m(g(t))\,dt\right)m_q(t_1).
\]
Two values belonging to different scales or schemes are not comparable
numbers until they are transported to the same fiber. This equation closes the
mathematical type of the running mass: it is not an isolated absolute scalar, but
a covariantly transported section.

\section{CKM Mixing and Covariance of Quark Spectral Frames}

\subsection{Unitary transport of the operator \texorpdfstring{\(B_d\)}{B_d} in the spectral frame of the type above by means of \texorpdfstring{\(V_{\rm CKM}\)}{V_CKM}}

The CKM mixing belongs to a distinct map. If \(B_u\) and \(B_d\) are the
mass operators in the up-type and down-type bases, the internal unitary matrix
satisfies
\[
 B_d^{(u)}=V_{\rm CKM} B_d V_{\rm CKM}^*.
\]

\subsection{Separation between the interior closure and the external statistical recognition}

The angles derived from \(A\), \(C^*\) and \(\Delta_4\), the phase
\(\delta_{\rm CP}\), the Jarlskog invariant and the determinant identity
of the commutator determine mixing without creating the mass eigenvalues. Thus
spectrum, color, generation and change of basis remain separate.

\section{Neutrino operator, PMNS mixing and Dirac--Majorana criteria}

\subsection{Three-dimensional neutral spectrum and inheritance of the action-clock-mass scale}

Neutrinos are electrically neutral fermionic sections. Their
fermionicity descends from the spinorial holonomy, the exterior completion, and
the canonical anticommutation relations; neutrality does not alter their
statistics. The internal support is
\[
 \mathcal H_\nu=
 \bigoplus_{j=1}^{4}\mathcal H_{\nu,G_j},
 \qquad
 \dim\mathcal H_\nu=2+4+6+8=20.
\]
The four structural depths are neither flavours nor eigenvalues. The
three-dimensional physical subspace is obtained through the neutral projector
\(P_\nu\), and the mass operator is
\[
 M_\nu=P_\nu\mathcal M_\nu P_\nu.
\]
Its three ordered eigenvalues are the masses; the absolute scale comes from the
same mass clock that fixes the remaining branches, and not from the mixing angles.

\subsection{Range 3 of the Gram matrix, polar frame, and antiunitary Majorana criterion}

Let \(M_\ell=P_\ell\mathcal M_\ell P_\ell\) be the charged operator. Denote by
\(F_\ell\) and \(F_\nu\) complete orthonormal spectral frames, ordered by
the spectral projectors of \(M_\ell\) and \(M_\nu\). The mixing matrix is
\[
 U_{\rm PMNS}=F_\ell^*F_\nu.
\]

\begin{theorem}[Spectral closure of the leptonic mixture]
If the two frames are complete, \(U_{\rm PMNS}\) is unitary. If the
spectra are simple, it is determined up to diagonal phase readjustments and
permutations fixed by the spectral order. If degeneracy exists, the
exact result is a double class
\[
 \prod_\lambda U(m_{\ell,\lambda})
 \backslash U(3) /
 \prod_\mu U(m_{\nu,\mu}),
\]
and not a numerically artificial matrix selected.
\end{theorem}

\begin{proof}
Orthonormality gives
\(U_{\rm PMNS}^*U_{\rm PMNS}=F_\nu^*F_\ell F_\ell^*F_\nu=I_3\).
Changes of basis within each degenerate subspace act on the
left or on the right, producing exactly the indicated
double class.
\end{proof}

The complete record determines the elements of \(M_\nu\) by means of carry,
sheet, orientation, vacancy, and holonomy. Equivalently, the
Hermitian kernel may be used
\[
 K_{\rm reg}(c,d)
 =\sum_{r\in\{\rm car,hoj,ori,vac,hol\}}
 \frac{\langle\Xi_r(c),\Xi_r(d)\rangle}
      {\sum_x\|\Xi_r(x)\|^2},
\]
with omission of every summand of null norm. Its Gram matrix is positive and
constructs the operator before diagonalization. The rank-one phase
\[
 I+(e^{i\pi/6}-1)|b_K\rangle\langle b_K|
\]
remains as a negative control: its formal unitarity is not sufficient and its angle
\(\theta_{13}\) falsifies its identification with full-rank physical
mixing.

Let \(G_0\) be the three-dimensional matrix obtained by aggregating the preceding kernel
over the three neutrino sectors. The finite condition
\[
 \det G_0\ne0
\]
is the exact certificate that the observables included in the record
separate three directions. In that case, the polar factor
\[
 U_{G_0}=G_0(G_0^*G_0)^{-1/2}
\]
is unitary and has rank \(3\). If \(\det G_0=0\), the declared kernel does not
determine a three-dimensional transport: an independent internal observable
must be added and the Gram matrix recomputed. Invertibility may not be
imposed by adding a term constructed from the mixing matrix one seeks to obtain.
The physical matrix is fixed by the spectral frames of \(M_\ell\) and
\(M_\nu\); the Gram criterion certifies sufficiency of the register, but does not
by itself select those frames.

Let \(\mathcal C_\nu\) be the antiunitary self-conjugation of the neutral sector. The
Majorana criterion is
\[
 \mathcal C_\nu^2=I,
 \qquad
 \mathcal C_\nu M_\nu=M_\nu\mathcal C_\nu,
 \qquad
 \mathcal C_\nu H_\nu=H_\nu\mathcal C_\nu,
\]
together with the existence of a fixed real basis in the realification of the
physical subspace. When these conditions are satisfied, each mode can be
represented by \(\psi=\mathcal C_\nu\psi\). The Dirac branch requires, by
contrast, a conserved charge \(Q\) such that
\[
 \mathcal C_\nu Q\mathcal C_\nu^{-1}=-Q,
 \qquad Q\psi\ne0,
\]
and retains \(\psi\) and \(\mathcal C_\nu\psi\) as distinct modes. If neither
set of conditions is established, the type remains undecided.
Neutrality, the word involution or the Ising sector, taken
separately, do not decide this alternative.

\section{Sterile neutrino sector and separation between existence and mixing}

\subsection{Orthogonal projector onto the active sector and cancellation of the weak current}

An additional structural depth does not by itself constitute a sterile
neutrino. A candidate sterile sector is defined by a nonzero projector
\(P_s\) satisfying
\[
 P_s P_{\rm act}=0,
 \qquad P_s J_{\text{weak}}=0,
\]
with dynamics defined on \(P_s\mathcal H_\nu\).

\subsection{Distinction between invariant reduction and non-vanishing coupling with the active sector}

If, in addition, \([P_s,M_\nu]=0\), the sector is reducing and does not mix with the active sector. If
\(P_sM_\nu P_{\rm act}\ne0\), there is active--sterile mixing and necessarily
\([P_s,M_\nu]\ne0\). Thus, the existence of the sector and the existence of
mixing are distinct properties. If the weakly neutral projector does not exist,
the additional coarse channel is an internal depth of the support and not a fourth
species.

\section{Bosons \texorpdfstring{\(W\) and \(Z\)}{W and Z} and scalar sector: projectors, resolvents, and pole criterion}

\subsection{Sectorial effective operator and meromorphic continuation of the reduced resolvent}

The sectors \(W\), \(Z\), and the scalar sector have their own
representation, route, and signature descriptors. Their realization does not require forcing them into an
elementary rank-one fiber. The positive fiber operator is constructed first
and then, on a finite-dimensional subspace \(P_X\mathcal H_X\), the
effective channel operator
\[
 H_{X,\rm eff}(z)
 =P_X H_X P_X
 +P_X H_X Q_X(z-Q_X H_X Q_X)^{-1}Q_X H_X P_X.
\]
\subsection{Isolated admissible zero, physical residue, and distinction with respect to mere channel opening}

If the reduced resolvent admits meromorphic continuation through the cut of the
continuum, the poles are the isolated zeros, counted with multiplicity, of
\[
 \det\bigl(z-H_{X,\rm eff}(z)\bigr)=0,
 \qquad
 z_X=M_X c_{\rm int}^2-\frac{i}{2}\Gamma_X.
\]
An open channel does not by itself guarantee the existence of such zeros. When
a zero with the prescribed physical residue exists, its real part publishes pole
mass and its imaginary part the width. For the
scalar sector, identification of an excitation additionally requires the scalar
projector and its gauge and fermionic couplings. Two independent poles
determine two states; a second label without projector or pole does not create a
particle.

\section{Composite states and cocyclic linking law}

\subsection{Single Meson and Baryon Projectors and Refinement of the Composite Signature}

For constituents \(X_1,\ldots,X_n\), the input is the composite record
\[
 \mathfrak C_{12}^{(n)}:
 (\Gamma_1,\ldots,\Gamma_n;\mathcal T,\partial)
 \longmapsto\mathcal L_{\rm comp},
\]
not the sum of isolated masses. The binary composition
\[
 L_1\star L_2=L_1+L_2+\mathfrak b(L_1,L_2)
\]
is associative exactly when
\[
 \mathfrak b(L_1,L_2)+\mathfrak b(L_1\star L_2,L_3)
 =\mathfrak b(L_2,L_3)+\mathfrak b(L_1,L_2\star L_3).
\]
The signature, the refinement, and the composite operator are computed after the linking:
\[
 \mathcal L_{\rm comp}\longmapsto\sigma_{\rm comp}
 \longmapsto\eta_{\rm comp}\longmapsto\mathcal M_{\rm comp}.
\]
The physical projector combines colour singlet, spin and parity, isospin, charge,
exchange symmetry, and channel.

\subsection{Feshbach--Schur Reduction and Construction of the Effective Operator}

The stable mass is an eigenvalue of the
compression; when the hypotheses for meromorphic continuation hold, the mass
and width of a resonance follow from a Feshbach pole. Consequently,
mesons, baryons, exotic states, nuclei, atoms, and
molecules share the same causal scheme without being reduced to a sum of masses.

Two constituent trees represent the same species if there exists a
unitary that simultaneously intertwines their operators, channel projectors and
poles. In the absence of such a unitary they are retained as distinct realizations,
even if they share constituents and partial quantum numbers.

\section{Resonances, widths, and survival of material sections}

\subsection{Complex Poles, Admissible Residues and Associated Widths for the Meromorphic Continuation}

Let \(S_X\) be the projector onto the surviving subspace and
\(T_X=S_X U S_X\) the compressed one-step operator. If it has a simple dominant
eigenvalue
\[
 \lambda_X=r_X e^{-i\theta_X},\qquad 0<r_X<1,
\]
fix the branch \(\theta_X\in I_\theta\), where \(I_\theta\) is the fundamental
quasienergy interval determined by the temporal chart and its orientation.
Then
\[
 M_X=\frac{\hbar_{\rm int}}{t_0c_{\rm int}^2}\theta_X,
 \qquad
 \Gamma_X=-\frac{2\hbar_{\rm int}}{t_0}\log r_X,
 \qquad
 \tau_X=\frac{\hbar_{\rm int}}{\Gamma_X}.
\]
\subsection{Compatibility of the genealogical refinement with resonant persistence}

Without the choice of \(I_\theta\), the phase determines only \(\theta_X\) modulo
\(2\pi\), and the quasienergy mass is not a well-defined scalar.
The width is not a sixth coordinate of the mass signature: it arises from the modulus
of the open mode. Equality of mass and width between particle and
antiparticle also requires transporting the channels and the causal prescription, not merely
the parity of the signature.

\section{Multiparticle completion and Fock statistics}

\subsection{Symmetric and Exterior Fock Functors Induced by Holonomy Parity}

Let \(\mathsf{Hol}_{\pm}^{\leq1}\) be the monoidal category of fibers of a
particle, with contractive refinement morphisms, in particular unitary ones,
and holonomy parity
\(\epsilon\in\{+1,-1\}\). The Fock completion is defined by the functor
\[
 \mathfrak F(H,\epsilon)=
 \begin{cases}
 \displaystyle\bigoplus_{n\ge0}\operatorname{Sym}^n(H),
     &\epsilon=+1,\\[2mm]
 \displaystyle\bigoplus_{n\ge0}\bigwedge^n\!H,
     &\epsilon=-1,
 \end{cases}
\]
and, for a contractive intertwiner \(T:H\to H'\),
\[
 \mathfrak F(T)=
 \bigoplus_{n\ge0}\operatorname{Sym}^n(T)
 \quad\text{or}\quad
 \bigoplus_{n\ge0}\bigwedge^n\!T.
\]

\subsection{CCR/CAR Relations, Pauli Exclusion and the Domain of Non-Contractive Morphisms}

\begin{theorem}[Internal transport of the statistic]
The foregoing completion is a functor of bounded operators on \(\mathsf{Hol}_{\pm}^{\leq1}\); it is compatible with refinement and is graded monoidal. In the \(\epsilon=-1\) branch, the creation and annihilation operators satisfy CAR and the occupation operators satisfy \(N_s^2=N_s\); in the \(\epsilon=+1\) branch, they satisfy CCR. Therefore, Pauli exclusion and fermionic and bosonic statistics descend from holonomy parity once the representation of a particle has been fixed.
\end{theorem}

\begin{proof}
Symmetric and exterior powers preserve compositions and identities; the contractivity of \(T\) ensures that their direct sum defines a bounded operator. For noncontractive morphisms, second quantization requires an explicit dense domain and does not belong to this category of bounded operators. The sign rule of the exterior product produces CAR and idempotence of occupation; the symmetric quotient produces CCR. Second quantization of the refinement morphisms intertwines the corresponding operators.
\end{proof}

This proof closes the internal algebraic transport. Identification
with the relativistic spin--statistics theorem additionally requires the
Lorentzian bridge and locality declared in Dimensional Mechanics; the two statements
must not be confused.

\section{Thermodynamic realization through Gibbs states and the pressure limit}

\subsection{Finite Hamiltonian, partition function, Gibbs state, and cofinal pressure limit}

For a compressed Hamiltonian \(H_{X,V}\) over a finite volume or refinement
\(V\), the thermal state is completely defined by
\[
 Z_{X,V}(\beta)=\operatorname{Tr}e^{-\beta H_{X,V}},
 \qquad
 \rho_{X,V,\beta}=Z_{X,V}(\beta)^{-1}e^{-\beta H_{X,V}}.
\]
The thermodynamic limit exists when the pressure
\[
 p_X(\beta)=\lim_{V\nearrow\infty}
 \frac1{\beta |V|}\log Z_{X,V}(\beta)
\]
exists and is independent of the cofinal sequence.

\subsection{Condensation criterion by saturation of the excited-state density}

Bosonic condensation is
determined by saturation of the excited-state density; it is not asserted
merely from the existence of a symmetric completion. The problem
is thus expressed by a verifiable mathematical condition on the
Hamiltonian and the limit, rather than by a nominal attribution.

\section{Topological sectors, virtual excitations, and classification of their codomains}

\subsection{Braided category, data \texorpdfstring{\(N\)--\(F\)--\(R\)}{N--F--R}, and pentagon and hexagon equations}

In a topological sector the primary codomain is a braided category, not
a line of masses. For simple objects \(a,b,c\), the physical data include
\[
 N_{ab}^{c},\qquad F^{abc}_{d},\qquad R^{ab}_{c},
\]
with pentagon and hexagon equations. A scalar mass exists only
when a Hamiltonian realization provides a spectral projector and a
gap. The Fibonacci, Ising/Majorana, and \(\mathbb Z_6\) sectors are thereby
closed as categorical codomains; they are not assigned a universal figure
by analogy with an elementary particle.

In the sector \(\mathbb Z_6\), fusion and braiding are
\[
 a\otimes b=a+b\pmod 6,
 \qquad
 R_{ab}=\zeta_6^{ab},
 \qquad \zeta_6=e^{2\pi i/6}.
\]
For Fibonacci,
\[
 \tau\otimes\tau=\mathbf 1\oplus\tau,
 \qquad
 F_\tau^{\tau\tau\tau}=
 \begin{pmatrix}
 \varphi^{-1}&\varphi^{-1/2}\\
 \varphi^{-1/2}&-\varphi^{-1}
 \end{pmatrix},
\]
\[
 R_{\mathbf1}^{\tau\tau}=e^{-4\pi i/5},
 \qquad
 R_\tau^{\tau\tau}=e^{3\pi i/5}.
\]
The Ising sector has objects \(\mathbf1,\psi,\sigma\) and rules
\[
 \psi\otimes\psi=\mathbf1,
 \qquad
 \psi\otimes\sigma=\sigma,
 \qquad
 \sigma\otimes\sigma=\mathbf1\oplus\psi,
\]
\[
 F_\sigma^{\sigma\sigma\sigma}
 =\frac1{\sqrt2}
 \begin{pmatrix}1&1\\1&-1\end{pmatrix},
 \qquad
 R_{\mathbf1}^{\sigma\sigma}=e^{-\pi i/8},
 \qquad
 R_\psi^{\sigma\sigma}=e^{3\pi i/8}.
\]
These matrices satisfy the coherence relations in the indicated
conventions. The duplication
\(\mathcal C\boxtimes\mathcal C^{\rm rev}\) preserves both orientations and
avoids introducing a chirality by external choice.

\subsection{Finite persistence of virtual states and absence of a universal scalar mass}

A virtual state is a section of finite persistence in an amplitude or a
resolvent. It does not belong to the asymptotic spectrum and does not require a real pole
mass. Its parameters are virtuality, the channel and the duration of
persistence. This definition avoids turning an internal computational line into
a new species.

\section{Gravity, torsion, and spectral criterion of the spin-2 collective mode}

\subsection{Tachyonic Exclusion by Positivity and Separation of Primary Gravity by Connection and Torsion}

Positivity excludes an elementary species with \(m^2<0\) within the
mass codomain. A negative eigenvalue of the Hessian of a background does not define a
tachyon as a particle: it diagnoses the instability of the background and requires changing
the solution around which linearization is performed.

Gravitation enters Dimensional Mechanics by means of connection, curvature, and
torsion. It does not require an elementary spin-\(2\) particle as a datum of
departure.

\subsection{Physical pole of the traceless transverse Hessian and positivity of the residue}

Let \(\mathcal S_{\rm grav}\) be the geometric action and
\(\mathcal H_{\rm TT}\) the transverse traceless subspace of its Hessian
at a solution. The linearized observable is
\[
 \mathcal K_{\rm TT}
 =P_{\rm TT}\,\delta^2\mathcal S_{\rm grav}\,P_{\rm TT}.
\]
A collective spin \(2\) mode exists as an asymptotic particle only
if the resolvent of \(\mathcal K_{\rm TT}\) has a physical pole with positive
residue and satisfies the constraints. If no such pole exists, the gravitational
interaction remains wholly in the connection and torsion. Thus, the
graviton is not a necessary component of the mass law, and its possible appearance
is decided by a spectral criterion, not by a nominal box.

\section{Equivalence between persistent section, inverse limit and finite particle representation}

\subsection{Quotient by tail equivalence and compatibility of section families}

The genealogical definition of a particle and its operator representation are
two presentations of the same object. Let
\[
 \mathfrak r_K=(x_K,\gamma_K,B_K,s_K,n_K,U_K)
\]
a finite representative at depth \(K\), with state, route, boundary,
orientation, counter, and representation, and let
\(r_{L,K}:\mathfrak R_L\to\mathfrak R_K\) be the refinement maps. A
persistent section is a family
\[
 X=(\mathfrak r_K)_{K\ge K_0},
 \qquad r_{L,K}(\mathfrak r_L)=\mathfrak r_K
 \quad(L\ge K\ge K_0).
\]
Two families are equivalent if they coincide on a cofinal tail.

\subsection{Unique prolongation condition to reconstruct a section from a finite evaluation}

\begin{theorem}[Equivalence of the Persistent Section and the Inverse Limit]
The category of persistent sections, defined as compatible families,
is canonically equivalent to
\[
 \varprojlim_K\mathfrak R_K/\!\sim_{\rm cola}.
\]
Projection to a finite depth is an evaluation, not an inverse functor. A representative \(\mathfrak r_K\) determines a unique persistent section if and only if it has a unique compatible prolongation modulo tail equivalence; this uniqueness is obtained, for example, when the refinement maps are bijective on the surviving tail from \(K\) onward.
\end{theorem}

\begin{proof}
By definition, a persistent section is precisely a point of the inverse
limit; the quotient identifies families that coincide on a cofinal tail.
Evaluation is the canonical projection onto one component. Its fiber contains
all compatible prolongations of the evaluated representative, so that the fiber
is a singleton exactly under the indicated unique-prolongation condition.
\end{proof}

An HMT--MD particle is therefore a class of persistent sections without a finite terminal obstruction, endowed with representation, orientation, charge, spin, and spectral character compatible with refinement. The route is its observable history, not its identity reduced to a cell. The antiparticle is the section conjugated by the dodecaphasic involution and the conjugate representation. An obstructed finite section is virtual; a cocyclic product of sections is composite; a persistent braided representation is anyonic.

\section{Equivariant realization closure theorem and exhaustive classification of material species}

\subsection{Comprehensive atlas of 324 routes: 49 invariants per section and conservation of the null sector}

\begin{theorem}[Closure by physical codomain]
Suppose that the atlas \(81/56/13/324\), the null sector, the
action--clock--mass scale, the readers \(K_D,K_\Omega\), the physical descriptor, and the
internal representation of each state have been constructed. Then every admissible and
realizable descriptor, that is, every \(X\) with
\(d_{\mathcal P}(X)\in\operatorname{im}d_{\mathcal R}\),
determines, without consulting its external value:
\begin{enumerate}
\item a fiber, orbit, or multisection \(\mathscr S(X)\);
\item the positive bidirectional pair and its joint PVM
      \((\widehat M_X^{\beta,D},\widehat M_X^{\beta,\Omega},
      E_X^{D,\Omega})\);
\item a decomposition into irreducible blocks;
\item a spectral measure for each physical coupling;
\item a scalar in a compressed rank-one sector or a multiplet in a
      reducible sector; in an open sector, an effective operator and, only if
      a meromorphic continuation with an admissible isolated zero exists, a
      complex pole;
\item a categorical object when the realization is topological and does not yet have
      a Hamiltonian gap;
\item a nonexistence criterion when the requested type contradicts
      positivity, gauge, or asymptoticity.
\end{enumerate}
None of these results requires selecting a route through mass
that will later be tested.
\end{theorem}

\begin{proof}
The fiber product constructs the fiber prospectively. The character and the
two readers construct the positive pair and its joint PVM. When the physical
type declares a scalar functional, functional calculus produces the corresponding
operator; conditional expectation projects it onto the commutant of the
symmetry. The isotypic decomposition produces scalars or blocks by Schur's
lemma. Preparation projectors produce spectral measures.
Feshbach reduction produces the effective operator; meromorphic continuation
and an admissible isolated zero produce a pole, if one exists. Topological
sectors retain their fusion and braiding data until a Hamiltonian
realization exists. Positivity and gauge constraints
exclude incompatible types. Each step depends only on internal data or
the declared coupling, never on the external column.
\end{proof}

\paragraph{Refutation controls.}

The construction is refuted in its domain if any of the following events occurs:

\begin{enumerate}
\item modifying the external values changes the fiber \(\mathscr S(X)\);
\item the assignment ceases to be equivariant under cellular inversion,
      conjugation, or gauge;
\item a point cell is selected in a noncentral multisection without a
      coupling that breaks the symmetry;
\item the synthesis of \(K_D\) and \(K_\Omega\) privileges an orientation or does not
      recover the diagonal when both readers coincide;
\item a compression with several eigenvalues is replaced by a mean without
      a physical projector;
\item the color quotient produces different masses for red, green, and blue;
\item quark masses are compared without declaring the same scheme and the same scale;
\item a rank-one phase is presented as PMNS mixing after failing the
      full-rank angular contrast;
\item neutrality is used by itself to assert the Majorana character;
\item the mass of a composite is obtained by adding constituent masses without
      a binding record;
\item the width is inserted as a coordinate of the signature instead of being extracted from
      an open mode or a pole;
\item a negative eigenvalue of the Hessian is published as a tachyonic particle;
\item gravitation is made dependent on an elementary graviton introduced
      before connection, curvature and torsion.
\end{enumerate}

\subsection{Complete inventories of 471 masses and 384 widths and subsequent metrological recognition}

External comparison begins only after the descriptor, fiber, operator, projector, scheme, and scale have been frozen. Each row must distinguish stable mass, pole mass, resonance estimator, and width. The atlas of \(324\) records and the inventories of \(471\) masses and \(384\) widths constitute the subsequent exhaustive presentation. The conceptual construction and each record cross-reference one another through their route identifiers. The records retain all their fields and are composed as coherent units; they are not replaced by abbreviated historical tables.

The comparison does not consist in enlarging a list of names, but in applying,
for each physical type, the complete map that leads from the genealogy of route
to its observable codomain.
